\section{Đối tượng khác nhau, cái gì đi qua được}
\label{sec:ch09-doi-tuong-khac-nhau}

\index{độ trung thực!phạm vi của bằng chứng}%
\index{chuỗi suy luận!so với attribution đặc trưng}%
Mọi con số của mục~\ref{sec:ch09-ket-qua} đều là điểm của một chỉ số chấm chuỗi
suy luận. Không chỉ số nào của chương~\ref{ch:do-do-trung-thuc} có mặt trong
bảng~\ref{tab:ch09-auroc}: deletion, insertion, comprehensiveness, sufficiency,
ROAR, Sensitivity-n và infidelity không xuất hiện ở đâu trong bài báo, và tám
chỉ số nó chấm đều được thiết kế riêng cho chuỗi suy luận. Vì vậy câu
\enquote{chỉ số trung thực không đo độ trung thực} đã được chứng minh bằng thực
nghiệm cho một lớp chỉ số, và lớp ấy không phải lớp mà Phần~II và
chương~\ref{ch:do-do-trung-thuc} đã dựng lại.

Ranh giới đó là ranh giới của chính bài báo chứ không phải một dè dặt sách này
thêm vào. Mục~\ref{sec:ch09-dinh-nghia-lai} đã dẫn chỗ các tác giả tự tách khỏi
lớp đối tượng của Phần~II, ngay tại chỗ họ viết định nghĩa
mới~\autocite{p21metrics}. Đọc kết quả của bài báo như một bản án dành cho LIME
hay SHAP là đi ngược lại lập luận mở đường cho chính các kết quả ấy.

Cái đi qua được ranh giới là ba thứ, và cả ba đều không phải con số.

Thứ nhất là phương pháp. Vòng luẩn quẩn mà
mục~\ref{sec:ch08-dieu-chua-giai-quyet} mô tả không phải đặc sản của chuỗi suy
luận; nó xuất hiện ở bất cứ đâu mà ground truth đòi biết nội dung tính toán bên
trong mô hình. Lối ra của bài báo cũng không gắn với chuỗi suy luận: chọn
những bài toán mà đầu ra để lộ phép tính đã sinh ra nó thì được ground truth
mà không cần nhìn vào trong. Nguyên tắc ấy chuyển sang attribution đặc trưng
được, và chương~\ref{ch:ground-truth-dung-san} nhận nó làm một trong ba lối
thoát của Phần~V.

Thứ hai là hai chẩn đoán, vì cả hai đều phát biểu được ở mức cấu trúc. Phép
đồng nhất mức quan trọng với độ trung thực không phải một sai lầm riêng của
Adding Mistakes; nó là chính phép suy luận mà cả họ xóa dần của
mục~\ref{sec:ch08-ho-xoa-dan} đứng trên, ở đó bỏ một
\tn{đặc trưng}{feature} mà đầu ra không
đổi được đọc thành mô hình không dùng đặc trưng ấy. Tương tự, việc một lời
biện minh bịa đặt vẫn qua được các chỉ số hữu dụng ngữ nghĩa là tính hợp lý
lọt vào dụng cụ đo, mà đó là rủi ro của mọi cách đánh giá lấy phản ứng của một
người đọc, hoặc một mô hình đọc, làm thước.

Thứ ba là một phép kiểm mà phía attribution có thể chạy ngay và phần lớn chưa
chạy: đo mức đồng thuận giữa các chỉ số. Kết quả kappa quanh 0 không cần ground
truth mới đọc được, mà chỉ cần chạy nhiều chỉ số trên cùng một tập. Nếu hai
chỉ số cùng đo một đại lượng thì chúng phải đồng thuận cao hơn thế nhiều, nên
mức đồng thuận thấp một mình nó đã đủ để loại bỏ khả năng cả hai cùng đo đúng
một đại lượng.

Cái không đi qua được thì gọn hơn: mọi con số. Bảng~\ref{tab:ch09-auroc} không
nói gì về việc deletion hay Sensitivity-n sẽ đạt AUROC bao nhiêu khi bị đem ra
đối chiếu tương tự, vì phép đối chiếu ấy chưa ai làm. Nói cho đúng thì trạng
thái hiện tại là thế này: với chuỗi suy luận, câu hỏi các chỉ số có đo đúng
thứ chúng ghi trên mặt đồng hồ hay không đã có câu trả lời và câu trả lời là
không; với attribution đặc trưng, câu hỏi ấy vẫn chưa được hỏi bằng thực
nghiệm. Đó là một khoảng trống, không phải một kết quả, và
chương~\ref{ch:khoang-trong-nghien-cuu} nhận nó về.

Ngay trong phạm vi của mình bài báo cũng để lại phần chưa xong, và các tác giả
ghi ra. Phương pháp chỉ cấp ground truth cho những bước mà thiết kế nhiệm vụ
buộc phải xảy ra, trong khi một chuỗi bất kỳ còn chứa nhiều bước khác mà thiết
kế ấy không nói được gì; thiết kế đặt độ chính xác lên trước cũng làm lệch
phân bố nhãn ở mức cả chuỗi theo hướng đã nói ở
mục~\ref{sec:ch09-bonafide}~\autocite{p21metrics}. Phương pháp cũng chỉ chạy
được trên mô hình trọng số mở, nên đúng những mô hình đóng đang được dùng
nhiều nhất lại nằm ngoài tầm. Và một ô của bảng bỏ trống vì chi phí, nghĩa là
có một chỉ số mà bộ đánh giá không đủ sức kiểm.

Điều đó đặt Phần~IV vào thế phải đi tiếp bằng đường khác. Nếu chưa ai đem các
chỉ số attribution ra đối chiếu với ground truth dựng sẵn, thì bằng chứng gần
nhất về tình trạng không trung thực phải tìm ở chỗ khác: ở những trường hợp mà
chính lời giải thích tự mâu thuẫn với nó. Chạy hai lần một phương pháp trên
cùng một mô hình và cùng một đầu vào mà ra hai lời giải thích khác nhau thì
không cần ground truth nào mới kết luận được rằng ít nhất một trong hai là
sai. Chương~\ref{ch:lime-nao-dang-tin} đọc trường hợp ấy trên chính LIME.
